% Section owners: all members (answers and limitations); Puneet (assembly)
\section{Discussion and Limitations}
\label{sec:p2discussion}
\textbf{Answers to the four questions.} (1)~The deeper residual network separates the confusable classes better:
\TestGainFinalVsPone{}~pp accuracy, Shirt recall \TestPoneShirt{} to \TestFinalShirt. The roadmap attributes the
gain to the architecture; the stronger recipe alone made the Phase 1 CNN slightly worse. (2)~Width, batch
normalisation, augmentation as a whole and the cosine schedule each account for about 0.7 to 1.1~pp; label
smoothing and random erasing for 0.2 to 0.3~pp; the horizontal flip, dropout and weight decay have no measurable
effect in this setting. (3)~Handcrafted and learned models do fail on different images, but feature-level fusion
recovers little of this: \HtwoGain{}~pp, consistent in sign and not significant. The practical value of the hybrid
work is the cascade, which matches the fusion accuracy while computing handcrafted features for \RoutedTest\% of
the images, and the observation that HOG alone is the most robust model under brightness and contrast changes.
(4)~The architecture gain is far outside both kinds of noise. It costs seven times the GPU inference time, worse
calibration and lower robustness to noise.

\textbf{Limitations.}
\begin{enumerate}
\item The test set was already used at the Phase 1 gate (Section~\ref{sec:p2protocol}).
\item Three seeds give a spread, not a significance test; McNemar tests use one seed.
\item The Phase 2 recipe was fixed in advance and not tuned. The ablations vary factors of the base model, not of
the deeper model that was finally selected; in particular, the shortcut ablation at seven layers says nothing about
thirteen.
\item One-factor ablations do not measure interactions, for example between dropout and augmentation.
\item Hybrid hyper-parameters were tuned on a 10{,}000-image subset where differences were within noise; the
selected $\alpha$ and number of HOG components lie at the edge of the grid; the SVM is fitted on embeddings of
images the CNN was trained on.
\item The robustness study uses 2{,}000 images and synthetic corruptions. It is an indication, not evidence about
photographs of real inventory.
\item Timings come from one shared cloud runtime. The run started from an uploaded archive, so the freeze file
records a timestamp and a hash but no git commit.
\item Fashion-MNIST contains ambiguous Shirt-group labels, which bounds the reachable accuracy.
\end{enumerate}

\textbf{Future work.} Temperature scaling on the validation set to restore calibration; noise and brightness
augmentation, followed by a repeat of the robustness test; out-of-fold embeddings for the hybrid; routing by input
statistics instead of confidence, so that images with unusual brightness reach the handcrafted model; and a test on
photographs taken outside the benchmark.

\section{Conclusion}
\label{sec:p2conclusion}
A thirteen-layer residual network with batch normalisation, global average pooling, augmentation and a cosine
schedule reaches \TestFinal{} $\pm$ \TestFinalStd{} test accuracy on Fashion-MNIST, \TestGainFinalVsPone{}~pp above
the Phase 1 CNN re-trained under the same protocol, with most of the gain in the upper-body classes. Controlled
ablations with three seeds identify which components matter and show that the training recipe helps only a model
with enough capacity. The handcrafted--learned fusion gives a small, consistent and statistically unsupported gain;
a confidence-gated cascade delivers the same accuracy at a fraction of its cost. The better model is less
calibrated and less robust to noise than its predecessor, which matters more for real inventory images than the
benchmark accuracy does.

\section*{Author Contributions}
M.~Selukar: architecture family, training engine, pipeline integration, experiments on the GPU.
A.~Singh: hybrid models, controls and cascade.
A.~Sharma: augmentation pipeline and robustness study.
T.~Rao: ablation design, regularisation analysis and model selection stage.
R.~Kumar: multi-seed statistics, significance tests, calibration and error analysis.
P.~Takhar: experiment configuration, run store, freeze and test gate, reporting and literature review.
All authors wrote the sections on their components and reviewed the manuscript.
